\section{Addressing, Growth, and the Exact Engine}

The crystal is an infinite object. To compute on it we have to name its cells, count how it expands, and
build it in memory without error. These three demands turn out to be one demand seen from three angles. The
same Fibonacci and golden-ratio structure that grows the crystal also addresses it, and the same exact
integer arithmetic that addresses it also builds it.

This section develops the golden ratio that pervades
the substrate, the Fibonacci addressing that names every cell, the graph the crystal becomes once geometry
is stripped away, the integer ladder that fixes the whole construction without free parameters, the growth
series that counts the expansion exactly, and the modular engine that realizes the crystal to tens of
millions of cells with no drift.

\subsection{The golden ratio and the Fibonacci numbers}

The Fibonacci sequence begins $1, 1, 2, 3, 5, 8, 13, 21, \dots$, each term the sum of the two before it,
$F(n) = F(n-1) + F(n-2)$. The ratios of successive terms converge on a single constant, the golden ratio
$\phi \approx 1.6180339887$, the number satisfying
\begin{equation}
  \phi^2 = \phi + 1.
\end{equation}
Squaring $\phi$ is the same as adding one to it. Dividing through gives $\phi = 1 + 1/\phi$, so $\phi$ is one
plus its own reciprocal. Its continued fraction is the all-ones expansion $1 + 1/(1 + 1/(1 + \cdots))$, which
is the slowest-converging of all continued fractions.

In the precise sense of continued-fraction
approximation, $\phi$ is the most irrational number. No simple fraction ever gets a good grip on it, which is
why packing at golden-ratio angles never lines up with itself and never leaves large gaps.

The golden ratio is not bolted onto the theory. It arrives in the $\{5,3,4\}$ substrate from three
independent directions, all verified together (P50).
\begin{itemize}
  \item First, as the fixed point of the self-referential equation $x = 1 + 1/x$, reached by pure iteration.
  \item Second, as the limiting growth ratio of the Fibonacci recurrence, a property of a growth process rather than an equation.
  \item Third, as the plane geometry of the pentagon, whose diagonal over its side equals $2\cos(\pi/5) = \phi$.
\end{itemize}
Arithmetic, a growth recurrence, and
plane geometry each give $\phi$, and the three agree to within one part in a million. When a single constant
arrives by three roads that share no machinery, this is the signature of one unified object rather than a
coincidence.

The third road explains why $\phi$ must appear at all. The leading $5$ in $\{5,3,4\}$ says the
faces of every cell are regular pentagons, and the pentagon and $\phi$ are the same fact wearing two hats. To
choose a pentagonal substrate is already to choose $\phi$, whether or not one intends to.

\subsection{Fibonacci addressing and greedy routing}

For a signal to cross the crystal, every cell needs a name from which it can compute its own neighbors with
no global map. Margenstern's splitting method supplies exactly this \cite{margenstern2013}. Pick any cell as
the root. Its immediate neighbors form ring one, their further neighbors ring two, and so on.

Following each
cell one step back toward the root traces a spanning tree, a skeleton that touches every cell exactly once.
Hyperbolic geometry is already tree-like, since the number of cells at distance $r$ grows like a constant to
the power $r$, so the spanning tree reads off a structure the geometry already carries rather than imposing
one.

The splitting method decomposes a wedge of the tiling into a leading tile plus smaller copies of a few
standard region types, recursing forever. That recursion is the tree, and counting new tiles per ring yields
the Fibonacci sequence $1, 2, 3, 5, 8, 13, \dots$. Number the cells ring by ring as the tree is walked
outward, then rewrite each number in the Fibonacci base.

By Zeckendorf's theorem every positive integer is a
unique sum of non-adjacent Fibonacci numbers. Writing that sum as a string of digits over the place values
$1, 2, 3, 5, 8, 13, \dots$ gives a clean, collision-free address for every cell. The first twelve are $1$,
$10$, $100$, $101$, $1000$, $1001$, $1010$, $10000$, $10001$, $10010$, $10100$, $10101$, as Table~\ref{tab:zeckendorf}
lays out. No address ever has
two ones in a row, because no integer ever needs two adjacent Fibonacci numbers. The set of valid addresses
is therefore a regular language, recognizable by a fixed-size machine with no global memory. A tiny automaton
inside every cell can navigate an infinite crystal.

\begin{table}[h]
\centering
\begin{tabular}{ccc}
\toprule
Cell number & Fibonacci sum & Address \\
\midrule
1  & $1$         & $1$     \\
2  & $2$         & $10$    \\
3  & $3$         & $100$   \\
4  & $3 + 1$     & $101$   \\
5  & $5$         & $1000$  \\
6  & $5 + 1$     & $1001$  \\
7  & $5 + 2$     & $1010$  \\
8  & $8$         & $10000$ \\
9  & $8 + 1$     & $10001$ \\
10 & $8 + 2$     & $10010$ \\
11 & $8 + 3$     & $10100$ \\
12 & $8 + 3 + 1$ & $10101$ \\
\bottomrule
\end{tabular}
\caption{The Zeckendorf addresses of the first twelve cells. Each cell number is the unique sum of non-adjacent
Fibonacci numbers, and no address has two adjacent ones.}
\label{tab:zeckendorf}
\end{table}

The address is a turn-by-turn route from the root. To find a cell's parent, strip one digit. To find its
children, extend by one legal step. To route from cell A to cell B, compare the two addresses from the left,
walk up A's route to the shared common ancestor by dropping digits, then walk down B's route by appending its
remaining digits.

Consider sending a message from cell $12$ at address $10101$ to cell $9$ at address
$10001$. Comparing from the left, $1$ matches, $0$ matches, then cell $12$ has $1$ where cell $9$ has $0$, so
the paths split. The shared prefix is $10$, which is the address of cell $2$, their common ancestor. Walking
up from $12$ by dropping one digit at a time gives $10101 \to 1010$ (cell $7$) $\to 101$ (cell $4$) $\to 10$
(cell $2$). Walking down to $9$ by appending its remaining digits gives $10 \to 100$ (cell $3$) $\to 1000$
(cell $5$) $\to 10001$ (cell $9$). The full path is $12, 7, 4, 2, 3, 5, 9$, six hops, and every cell on it
only compared two short strings. Each address in this route can be read straight from Table~\ref{tab:zeckendorf}.

No cell ever consulted a map of the crystal. On the 2D heptagrid this exact
addressed routing delivered $100$ percent of random pairs at low stretch (P42). On the real 3D dodecagrid,
greedy hyperbolic routing, where each cell forwards to the neighbor whose address is closest to the target,
delivered over $90$ percent of pairs at stretch under $2$ (P76). The addresses are not stored labels. They
are implicit in each cell's place in the web and can be computed from local structure whenever needed.

\subsection{The crystal as a graph}

Strip away the drawing, the cell shapes, and the curved-space picture, and what remains is a graph. Each cell
is a vertex, each shared face an edge. A dodecahedron has twelve faces, so in the interior every cell touches
exactly twelve others. The graph is $12$-regular. Distance is step-counting, the fewest edges between two
cells, and this step-counting distance is what space will mean.

Because the crystal lives in negatively
curved space, the graph is an expander, with no bottlenecks. Any region not larger than half the graph has a
boundary that is a fixed fraction of its size, so the graph cannot be cut into two large pieces by snipping
few edges. This is the same exponential ball growth seen as connectivity. The graph is also tree-like, since
the exponential branching that makes it an expander is the exponential ball growth that makes it tree-like,
two views of one fact.

Together these force a tiny diameter, logarithmic in the cell count. On roughly
$120{,}000$ exact $\{5,3,4\}$ cells the worst-case separation is about six hops through the bulk, close to
$\log_B N$ with growth base $B \approx 7.87$ (P111). Among a hundred thousand cells, no two are more than
about six steps apart. That tiny $O(\log N)$ diameter is the engine behind fast global integration, against
the $N^{1/3}$ crawl a flat cubic lattice would require.

\subsection{The integer ladder}

The entire substrate is built from the plain whole numbers with no free parameters at any step. One picks the
integers, which cannot be picked differently, and the crystal $\{5,3,4\}$, which is itself forced. Everything
else follows by deterministic construction. The ladder runs upward through four rungs, each verified (P48,
P104, P51, P83).

Asking the integers for their symmetries returns the modular group $\mathrm{PSL}(2,\mathbb{Z})$,
the integer matrices of determinant one, with nothing to tune. This parameter-free modular base tessellates
hyperbolic space Lorentz-safely, is addressed by a deterministic finite automaton through continued
fractions, and carries the golden ratio as its central geodesic, the all-ones continued fraction whose
convergents are the Fibonacci ratios (P48).

The next rung represents the exact dodecahedral cells of
$\{5,3,4\}$ in pure integer arithmetic (P104). Above it the whole tower runs as one continuous program, from
integer generator data up to running physics, through a single shared pipeline for three different bases
(P51). The top rung grows the crystal one cell at a time by a fixed deterministic rule, resumable,
append-only, and faithful to the static tiling ring for ring, with the hyperbolic curvature and the golden
ratio emerging on their own rather than being inserted (P83).

The integers are the bones, the small ones
carry the laws, and there are essentially no free parameters anywhere in the foundation.

\subsection{The growth series}

How fast does the crystal expand? Stand on one cell and count the cells one step out, two steps out, three
steps out. The growth function $f(n)$ records the number of cells added at ring $n$. For the dodecagrid the
first counts are $1, 12, 102, 812, 6402, 50412, \dots$.

These counts are fixed exactly by the same mirror
data that defines the crystal. Steinberg's formula computes them by packing the sequence into a generating
function $W(t) = \sum_n f(n)\, t^n$ and expressing its reciprocal as an alternating sum over the finite
parabolic subgroups,
\begin{equation}
  \frac{1}{W(t)} = \sum_{\text{finite } T} \frac{(-1)^{|T|}}{f_T(t)},
\end{equation}
where each $f_T(t)$ is the Poincare growth polynomial of a finite reflection subgroup. The whole computation
runs in exact integer and rational arithmetic over polynomials, with no rounding.

The denominator of $W(t)$
gives a linear recurrence. For $\{5,3,4\}$ the recurrence is
\begin{equation}
  f(n) = 9\, f(n-1) - 9\, f(n-2) + f(n-3),
\end{equation}
seeded by $1, 12, 102, 812$. Verifying on the seeds, the next term is
$f(4) = 9(812) - 9(102) + 12 = 7308 - 918 + 12 = 6402$, and the one after is
$f(5) = 9(6402) - 9(812) + 102 = 50412$, both matching the series exactly. The denominator $1 - 9t + 9t^2 -
t^3$ has characteristic polynomial $x^3 - 9x^2 + 9x - 1 = (x-1)(x^2 - 8x + 1)$, with roots $1$, $4 +
\sqrt{15}$, and $4 - \sqrt{15}$. The growth base is
\begin{equation}
  B = 4 + \sqrt{15} \approx 7.873.
\end{equation}
The reciprocal root $4 - \sqrt{15} \approx 0.127$ lies strictly inside the unit circle, which is the proof of
exponential growth. Each ring is about $7.87$ times the size of the one before it, and that is the universe's
expansion rate in cells. A direct breadth-first count returns the same $1, 12, 102, 812, 6402, 50412$, with
the ring-to-ring ratio climbing $8.50, 7.96, 7.88, 7.87$ toward $4 + \sqrt{15}$.

Table~\ref{tab:layers} lists the first ten rings of the two-dimensional heptagrid $\{7,3\}$ and the
three-dimensional dodecagrid $\{5,3,4\}$ side by side. The heptagrid grows at the golden ratio squared,
about $2.618$ per ring, and the dodecagrid grows at $4 + \sqrt{15}$, about $7.87$ per ring. The dodecagrid
reaches more than a billion cells by the tenth ring.

\begin{table}[h]
\centering
\begin{tabular}{rrrrr}
\toprule
 & \multicolumn{2}{c}{Heptagrid $\{7,3\}$} & \multicolumn{2}{c}{Dodecagrid $\{5,3,4\}$} \\
\cmidrule(lr){2-3} \cmidrule(lr){4-5}
Ring & Added & Cumulative & Added & Cumulative \\
\midrule
0 & 1 & 1 & 1 & 1 \\
1 & 7 & 8 & 12 & 13 \\
2 & 21 & 29 & 102 & 115 \\
3 & 56 & 85 & 812 & 927 \\
4 & 147 & 232 & 6{,}402 & 7{,}329 \\
5 & 385 & 617 & 50{,}412 & 57{,}741 \\
6 & 1{,}008 & 1{,}625 & 396{,}902 & 454{,}643 \\
7 & 2{,}639 & 4{,}264 & 3{,}124{,}812 & 3{,}579{,}455 \\
8 & 6{,}909 & 11{,}173 & 24{,}601{,}602 & 28{,}181{,}057 \\
9 & 18{,}088 & 29{,}261 & 193{,}688{,}012 & 221{,}869{,}069 \\
10 & 47{,}355 & 76{,}616 & 1{,}524{,}902{,}502 & 1{,}746{,}771{,}571 \\
\bottomrule
\end{tabular}
\caption{Cells added and cumulative total per ring for the two crystals. The heptagrid grows by about
$2.618$ per ring (the golden ratio squared), the dodecagrid by about $7.87$ (the root $4 + \sqrt{15}$).}
\label{tab:layers}
\end{table}

Two distinct sequences must
be kept apart. The cell-shell growth above counts honeycomb cells at each distance from a seed cell. The
Coxeter group growth, which Steinberg's formula also produces, counts group elements by word length, giving
$1, 4, 9, 17, 29, 46, \dots$ for $\{5,3,4\}$, a different object. The cell-shell series is the physically
meaningful one.

The two-dimensional warm-ups $\{7,3\}$ and $\{5,4\}$ share the denominator $1 - 3t + t^2$ and
grow at base $(3 + \sqrt5)/2 \approx 2.618 = \phi^2$, exactly the square of the golden ratio, so the
golden thread runs through the smaller cases as well.

\subsection{The exact modular-fingerprint engine}

Building the crystal by float coordinates fails at a hard precision wall. Tracking each cell by its center as
$g \cdot c_0$ projected into the Poincare ball works only while distinct cell centers stay further apart than
a $64$-bit float can resolve. Hyperbolic space grows exponentially and cells crowd toward the boundary of the
ball, so their coordinates pile up with gaps shrinking exponentially with distance. In practice the float
engine caps out near $15{,}500$ cells. This is not a coding flaw but the geometry.

The escape is to notice
that the reflection matrices of $\{5,3,4\}$ are made of a small set of exact algebraic numbers, with the only
irrational ingredients being $\sqrt5$ inside $\phi$ and $\sqrt2$. A finite field of integers modulo a prime
$p$ contains both square roots exactly when $p \equiv 1 \pmod{40}$, since then $2$ and $5$ are both quadratic
residues. The roots are computed by Tonelli-Shanks, $\phi$ is built from $\sqrt5$, and every dodecahedral
reflection becomes exact integer arithmetic with zero rounding.

A single prime carries a tiny risk that two
distinct cells share one residue, so the engine uses two primes at once, near $67{,}108{,}837$ and near
$66{,}000{,}000$, both congruent to $1 \bmod 40$. A cell's fingerprint is its center reduced modulo both
primes, and two cells collide only if they collide modulo both at once, which by the Chinese Remainder
Theorem means agreeing modulo a product above $4 \times 10^{15}$. The twelve face reflections are derived by
conjugating the outer generator by each of the $120$ elements of the dodecahedral cell stabilizer, so
adjacency is read straight off the group rather than measured. A typed open-addressing hash on the packed
fingerprints deduplicates cells with no allocation.

The result has been run offline to $20$ million cells at
a few seconds per million, every cell carrying exactly twelve facet neighbors throughout, verified against an
independent float build below the wall and persisted to disk without loss (P104). For transport tests, which
a tiny-diameter ball cannot support, the same exact arithmetic builds a sliver, a long thin geodesic tube
with bounded width whose alternating face product is a hyperbolic translation, giving genuine distance to
travel along a one-dimensional spine. The physics of the theory is therefore measured on the true crystal,
not on a guess of it.
